\documentclass[a4paper]{article}

\usepackage{fancyhdr}
\usepackage{extramarks}
\usepackage{amsmath}
\usepackage{amsthm}
\usepackage{amsfonts}
\usepackage{amssymb}
\usepackage{tikz}
% \usepackage[plain]{algorithm}
% \usepackage{algpseudocode}
\usepackage{enumerate}

\usetikzlibrary{automata,positioning}

% =========================================================
% Basic Document Settings
% =========================================================

\topmargin=-0.2in
\evensidemargin=0in
\oddsidemargin=0in
\textwidth=6.5in
\textheight=9.5in
\headsep=0.25in

\linespread{1.1}
\allowdisplaybreaks

% =========================================================
% Homework Details
% =========================================================

\newcommand{\hmwkTitle}{Homework\ \#09}
\newcommand{\hmwkClass}{Probability \& Statistics for EECS}
\newcommand{\hmwkDueDate}{Nov. 29, 2026, 23:59 UTC+8}
\newcommand{\hmwkAuthorName}{Zhang Yiheng}
\newcommand{\hmwkAuthorID}{2025531068}

\pagestyle{fancy}
\lhead{\hmwkAuthorName}
\chead{\hmwkClass : \hmwkTitle}
\rhead{\firstxmark}
\lfoot{\lastxmark}
\cfoot{\thepage}

\renewcommand\headrulewidth{0.4pt}
\renewcommand\footrulewidth{0.4pt}
\setlength\parindent{0pt}

% =========================================================
% Problem Environment
% =========================================================

\newcommand{\enterProblemHeader}[1]{
    \nobreak
    \extramarks{}{
        Problem \arabic{#1} continued on next page\ldots
    }
    \nobreak{}
    \nobreak
    \extramarks{
        Problem \arabic{#1} (continued)
    }{
        Problem \arabic{#1} continued on next page\ldots
    }
    \nobreak{}
}

\newcommand{\exitProblemHeader}[1]{
    \nobreak
    \extramarks{
        Problem \arabic{#1} (continued)
    }{
        Problem \arabic{#1} continued on next page\ldots
    }
    \nobreak{}
    \stepcounter{#1}
    \nobreak
    \extramarks{Problem \arabic{#1}}{}
    \nobreak{}
}

\setcounter{secnumdepth}{0}

\newcounter{partCounter}
\newcounter{homeworkProblemCounter}
\setcounter{homeworkProblemCounter}{1}

\nobreak
\extramarks{
    Problem \arabic{homeworkProblemCounter}
}{}
\nobreak{}

\newenvironment{homeworkProblem}[2][-1]{%
    \ifnum#1>0\relax
        \setcounter{homeworkProblemCounter}{#1}
    \fi
    \section{
        Problem \arabic{homeworkProblemCounter}
        (#2 points)
    }
    \setcounter{partCounter}{1}
    \enterProblemHeader{homeworkProblemCounter}
}{%
    \exitProblemHeader{homeworkProblemCounter}
}

% =========================================================
% Useful Commands
% =========================================================

\newcommand{\E}{\mathrm{E}}
\newcommand{\Var}{\mathrm{Var}}
\newcommand{\Cov}{\mathrm{Cov}}
\newcommand{\Bias}{\mathrm{Bias}}
\newcommand{\dd}{\,\mathrm{d}}
\newcommand{\diff}{\mathrm{d}}

\newcommand{\solution}{\medskip\noindent\textbf{\large Solution}\medskip}

\author{
    Name: \textbf{\hmwkAuthorName}\\
    Student ID: \hmwkAuthorID
}

\date{}

% =========================================================
% Document
% =========================================================

\begin{document}

\title{
SI 140A.02 Probability \& Statistics for EECS,
Fall 2026\\
Homework 9
}

\date{Due on \hmwkDueDate}
\maketitle

Read all the instructions below carefully before you
start working on the assignment, and before you
make a submission.

\begin{itemize}
    \item You are required to submit your homework
    in PDF format generated by \textbf{LaTeX},
    otherwise you will get no points of this homework.

    \item You are required to write down all the
    major steps towards making your conclusions;
    otherwise you may obtain limited points
    of the problem.

    \item Write your homework in English;
    otherwise you will get no points of this homework.

    \item Any form of plagiarism will lead to
    $0$ point of this homework.
\end{itemize}

\newpage

% =========================================================
% Problem 1
% =========================================================

\begin{homeworkProblem}{15}

Let $X$ be a continuous random variable with PDF

\[
f_X(x)=
\begin{cases}
cx(2-x), & 0<x<2,\\
0, & \text{otherwise},
\end{cases}
\]

where $c$ is a positive constant.

\begin{enumerate}[(a)]

\item Determine $c$ and find the CDF $F_X(x)$
for all real $x$.

\item Compute
\[
P\left(X\geq\frac32\mid X\geq1\right)
\quad\text{and}\quad
\E[|X-1|].
\]

\end{enumerate}

\solution

\begin{enumerate}[(a)]

\item \textbf{Determining the normalization constant $c$.}
Since $f_X$ is a valid PDF, it must satisfy the normalization condition
\[
\int_{-\infty}^{\infty} f_X(x)\dd x = 1.
\]
Because $f_X(x) = 0$ outside $(0,2)$, this reduces to
\begin{align*}
1
&= \int_{0}^{2} c x(2-x)\dd x
 = c \int_{0}^{2} (2x - x^2)\dd x \\
&= c \left[ x^2 - \frac{x^3}{3} \right]_{0}^{2}
 = c \left( 2^2 - \frac{2^3}{3} \right) - c(0) \\
&= c \left( 4 - \frac{8}{3} \right)
 = \frac{4}{3} c.
\end{align*}
Solving for $c$ yields
\[
\boxed{c = \frac{3}{4}}\,.
\]
Note that $c = \frac{3}{4} > 0$, and $f_X(x) = \frac{3}{4}x(2-x) > 0$ for all $x \in (0,2)$, so $f_X$ is indeed a valid PDF.

\medskip
\textbf{Finding the CDF $F_X(x)$ for all real $x$.}
By definition, the CDF is given by
\[
F_X(x) = P(X \le x) = \int_{-\infty}^{x} f_X(t)\dd t.
\]
We consider three regions for $x \in \mathbb{R}$:
\begin{itemize}
    \item \emph{Case 1: $x \le 0$.} Since $f_X(t) = 0$ for all $t \le 0$,
    \[
    F_X(x) = \int_{-\infty}^{x} 0\dd t = 0.
    \]

    \item \emph{Case 2: $0 < x < 2$.} Splitting the integral at $t = 0$, we obtain
    \begin{align*}
    F_X(x)
    &= \int_{-\infty}^{0} 0\dd t + \int_{0}^{x} \frac{3}{4} t(2-t)\dd t
     = \frac{3}{4} \int_{0}^{x} (2t - t^2)\dd t \\
    &= \frac{3}{4} \left[ t^2 - \frac{t^3}{3} \right]_{0}^{x}
     = \frac{3}{4} \left( x^2 - \frac{x^3}{3} \right)
     = \frac{3}{4} x^2 - \frac{1}{4} x^3
     = \frac{x^2(3-x)}{4}.
    \end{align*}

    \item \emph{Case 3: $x \ge 2$.} Since the entire probability mass lies on $(0,2)$,
    \[
    F_X(x) = \int_{0}^{2} f_X(t)\dd t + \int_{2}^{x} 0\dd t = 1 + 0 = 1.
    \]
\end{itemize}
(As a consistency check, the middle branch equals $0$ at $x = 0$ and $1$ at $x = 2$, so $F_X$ is continuous on $\mathbb{R}$.) Combining all three cases,
\[
\boxed{
F_X(x) =
\begin{cases}
0, & x \le 0, \\[4pt]
\dfrac{3}{4}x^2 - \dfrac{1}{4}x^3, & 0 < x < 2, \\[8pt]
1, & x \ge 2.
\end{cases}
}
\]

\item \textbf{Computing $P\!\left(X \ge \frac{3}{2} \mid X \ge 1\right)$.}
Since $\left\{X \ge \frac{3}{2}\right\} \subseteq \{X \ge 1\}$, the definition of conditional probability gives
\[
P\left(X \ge \frac{3}{2} \;\middle|\; X \ge 1\right)
= \frac{P\left(X \ge \frac{3}{2},\, X \ge 1\right)}{P(X \ge 1)}
= \frac{P\left(X \ge \frac{3}{2}\right)}{P(X \ge 1)}.
\]
Because $X$ is a continuous random variable, $P(X \ge a) = 1 - F_X(a)$ for any $a \in \mathbb{R}$. Using the CDF from part (a):
\begin{align*}
F_X(1) &= \frac{3}{4}(1)^2 - \frac{1}{4}(1)^3 = \frac{3}{4} - \frac{1}{4} = \frac{1}{2}
\implies P(X \ge 1) = 1 - \frac{1}{2} = \frac{1}{2}, \\[4pt]
F_X\!\left(\frac{3}{2}\right) &= \frac{3}{4}\left(\frac{3}{2}\right)^2 - \frac{1}{4}\left(\frac{3}{2}\right)^3
= \frac{27}{16} - \frac{27}{32} = \frac{27}{32}
\implies P\!\left(X \ge \frac{3}{2}\right) = 1 - \frac{27}{32} = \frac{5}{32}.
\end{align*}
Therefore,
\[
P\left(X \ge \frac{3}{2} \;\middle|\; X \ge 1\right)
= \frac{5/32}{1/2}
= \boxed{\frac{5}{16}}\,.
\]

\medskip
\textbf{Computing $\E[|X-1|]$.}
By the Law of the Unconscious Statistician (LOTUS),
\[
\E[|X-1|] = \int_{-\infty}^{\infty} |x-1|\, f_X(x)\dd x
= \int_{0}^{2} |x-1| \cdot \frac{3}{4}x(2-x)\dd x.
\]
Make the substitution $u = x - 1$, so $x = 1 + u$, $2 - x = 1 - u$, $\diff x = \diff u$, and $x(2-x) = (1+u)(1-u) = 1 - u^2$. As $x$ ranges from $0$ to $2$, $u$ ranges from $-1$ to $1$. Since $|u|(1-u^2)$ is an even function of $u$ on $[-1,1]$, splitting at $u = 0$ yields two equal integrals:
\begin{align*}
\E[|X-1|]
&= \frac{3}{4} \int_{-1}^{1} |u|(1 - u^2)\dd u
 = \frac{3}{4} \left( \int_{-1}^{0} (-u)(1 - u^2)\dd u + \int_{0}^{1} u(1 - u^2)\dd u \right) \\
&= 2 \cdot \frac{3}{4} \int_{0}^{1} u(1 - u^2)\dd u
 = \frac{3}{2} \int_{0}^{1} (u - u^3)\dd u \\
&= \frac{3}{2} \left[ \frac{u^2}{2} - \frac{u^4}{4} \right]_{0}^{1}
 = \frac{3}{2} \left( \frac{1}{2} - \frac{1}{4} \right)
 = \frac{3}{2} \cdot \frac{1}{4}
 = \boxed{\frac{3}{8}}\,.
\end{align*}

\end{enumerate}

\end{homeworkProblem}

\newpage

% =========================================================
% Problem 2
% =========================================================

\begin{homeworkProblem}{15}

Consider the following CDF:

\[
F(x)=
\begin{cases}
0, & x\leq0,\\[2pt]
\dfrac{x^2}{2}, & 0<x\leq1,\\[6pt]
1-\dfrac{(2-x)^2}{2}, & 1<x<2,\\[6pt]
1, & x\geq2.
\end{cases}
\]

\begin{enumerate}[(a)]

\item Let $U\sim\mathrm{Unif}(0,1)$.
Use inverse transform sampling to express
$X=F^{-1}(U)$ as a function of $U$.

\item For a random variable $X$ with CDF $F$,
compute
\[
P\left(\frac12<X<\frac32\right)
\]
and use the survival function to find $\E[X]$.

\end{enumerate}

\solution

\begin{enumerate}[(a)]

\item Observe first that $F$ is continuous on $\mathbb{R}$ and strictly increasing on $[0,2]$, with boundary values
\[
F(0) = 0, \qquad
F(1) = \frac{1^2}{2} = \frac{1}{2} = \lim_{x \to 1^{+}} \left(1 - \frac{(2-x)^2}{2}\right), \qquad
F(2) = 1.
\]
Thus $F$ maps $(0,2)$ bijectively onto $(0,1)$, and we can invert $u = F(x)$ on the two subintervals of $(0,1)$:
\begin{itemize}
    \item \emph{Case 1: $0 < u \le \frac{1}{2}$ (corresponding to $0 < x \le 1$).}
    Setting $u = F(x)$ gives
    \[
    u = \frac{x^2}{2} \implies x^2 = 2u.
    \]
    Since $x > 0$, we take the positive square root:
    \[
    x = F^{-1}(u) = \sqrt{2u}.
    \]

    \item \emph{Case 2: $\frac{1}{2} < u < 1$ (corresponding to $1 < x < 2$).}
    Setting $u = F(x)$ gives
    \[
    u = 1 - \frac{(2-x)^2}{2} \implies \frac{(2-x)^2}{2} = 1 - u \implies (2-x)^2 = 2(1-u).
    \]
    Since $1 < x < 2$, we have $2 - x > 0$, so taking the positive square root gives $2 - x = \sqrt{2(1-u)}$, and hence
    \[
    x = F^{-1}(u) = 2 - \sqrt{2(1-u)}.
    \]
\end{itemize}
Therefore, for $U \sim \mathrm{Unif}(0,1)$ (which lies in $(0,1)$ with probability $1$), the inverse transform sample $X = F^{-1}(U)$ is
\[
\boxed{
X = F^{-1}(U) =
\begin{cases}
\sqrt{2U}, & 0 < U \le \dfrac{1}{2}, \\[8pt]
2 - \sqrt{2(1-U)}, & \dfrac{1}{2} < U < 1.
\end{cases}
}
\]
(Note that both branches agree at $U = \frac{1}{2}$, giving $X = 1$, and extend continuously to $F^{-1}(0) = 0$ and $F^{-1}(1) = 2$.)

\item \textbf{Computing $P\!\left(\frac{1}{2} < X < \frac{3}{2}\right)$.}
Because $F$ is continuous everywhere, $X$ is a continuous random variable and $P(X = a) = 0$ for every $a \in \mathbb{R}$. Consequently,
\[
P\left(\frac{1}{2} < X < \frac{3}{2}\right) = F\!\left(\frac{3}{2}\right) - F\!\left(\frac{1}{2}\right).
\]
Evaluating $F$ at the two endpoints:
\begin{align*}
F\!\left(\frac{1}{2}\right) &= \frac{(1/2)^2}{2} = \frac{1/4}{2} = \frac{1}{8}, \\[4pt]
F\!\left(\frac{3}{2}\right) &= 1 - \frac{(2 - 3/2)^2}{2} = 1 - \frac{(1/2)^2}{2} = 1 - \frac{1}{8} = \frac{7}{8}.
\end{align*}
Hence,
\[
P\left(\frac{1}{2} < X < \frac{3}{2}\right) = \frac{7}{8} - \frac{1}{8} = \frac{6}{8} = \boxed{\frac{3}{4}}\,.
\]

\medskip
\textbf{Finding $\E[X]$ using the survival function.}
The survival function of $X$ is $S(x) = P(X > x) = 1 - F(x)$, which equals
\[
S(x) =
\begin{cases}
1, & x \le 0, \\[4pt]
1 - \dfrac{x^2}{2}, & 0 < x \le 1, \\[8pt]
\dfrac{(2-x)^2}{2}, & 1 < x < 2, \\[8pt]
0, & x \ge 2.
\end{cases}
\]
Since $F(0) = 0$, we have $P(X \ge 0) = 1$ ($X$ is nonnegative). By the survival-function formula for the expectation of a nonnegative random variable,
\[
\E[X] = \int_{0}^{\infty} S(x)\dd x
= \int_{0}^{1} \left(1 - \frac{x^2}{2}\right)\dd x
+ \int_{1}^{2} \frac{(2-x)^2}{2}\dd x
+ \int_{2}^{\infty} 0\dd x.
\]
We evaluate each of the two nonzero integrals separately:
\begin{itemize}
    \item For the first integral,
    \[
    \int_{0}^{1} \left(1 - \frac{x^2}{2}\right)\dd x
    = \left[ x - \frac{x^3}{6} \right]_{0}^{1}
    = \left(1 - \frac{1}{6}\right) - 0
    = \frac{5}{6}.
    \]

    \item For the second integral, substitute $w = 2 - x$, so $\diff w = -\diff x$. When $x = 1$, $w = 1$; when $x = 2$, $w = 0$. Therefore,
    \[
    \int_{1}^{2} \frac{(2-x)^2}{2}\dd x
    = \int_{1}^{0} \frac{w^2}{2}\,(-\diff w)
    = \int_{0}^{1} \frac{w^2}{2}\dd w
    = \left[ \frac{w^3}{6} \right]_{0}^{1}
    = \frac{1}{6}.
    \]
\end{itemize}
Summing the two contributions yields
\[
\E[X] = \frac{5}{6} + \frac{1}{6} = \boxed{1}\,.
\]

\end{enumerate}

\end{homeworkProblem}

\newpage

% =========================================================
% Problem 3
% =========================================================

\begin{homeworkProblem}{15}

Let $X\sim\mathrm{Expo}(\lambda)$, where
$\lambda>0$, and let $U\sim\mathrm{Unif}(0,1)$.
Assume that $X$ and $U$ are independent.

Consider the event
\[
A=\{U<e^{-\lambda X}\}.
\]

Find $P(X\leq t\mid A)$ for all real $t$,
and identify the corresponding distribution.

\solution

Since $X \sim \mathrm{Expo}(\lambda)$ and $U \sim \mathrm{Unif}(0,1)$ are independent, their marginal PDFs are
\[
f_X(x) =
\begin{cases}
\lambda e^{-\lambda x}, & x > 0, \\
0, & x \le 0,
\end{cases}
\qquad\text{and}\qquad
f_U(u) =
\begin{cases}
1, & 0 < u < 1, \\
0, & \text{otherwise},
\end{cases}
\]
so their joint PDF is
\[
f_{X,U}(x,u) = f_X(x) f_U(u) =
\begin{cases}
\lambda e^{-\lambda x}, & x > 0 \text{ and } 0 < u < 1, \\
0, & \text{otherwise}.
\end{cases}
\]
Note that for every $x > 0$, since $\lambda > 0$, we have $0 < e^{-\lambda x} < 1$. Therefore, inside the support $(0,\infty) \times (0,1)$, the condition $u < e^{-\lambda x}$ holds precisely for $0 < u < e^{-\lambda x}$.

\medskip
\textbf{Step 1: Compute $P(A)$.}
Integrating the joint PDF over the region $\{(x,u) : x > 0,\ 0 < u < e^{-\lambda x}\}$, we obtain
\begin{align*}
P(A)
= P\!\left(U < e^{-\lambda X}\right)
&= \int_{0}^{\infty} \int_{0}^{e^{-\lambda x}} \lambda e^{-\lambda x} \dd u \dd x \\
&= \int_{0}^{\infty} \lambda e^{-\lambda x} \Big[ u \Big]_{0}^{e^{-\lambda x}} \dd x \\
&= \int_{0}^{\infty} \lambda e^{-\lambda x} \cdot e^{-\lambda x} \dd x \\
&= \int_{0}^{\infty} \lambda e^{-2\lambda x} \dd x.
\end{align*}
Evaluating this improper integral gives
\[
P(A)
= \lim_{b \to \infty} \int_{0}^{b} \lambda e^{-2\lambda x} \dd x
= \lim_{b \to \infty} \left[ -\frac{1}{2} e^{-2\lambda x} \right]_{0}^{b}
= \lim_{b \to \infty} \left( -\frac{1}{2} e^{-2\lambda b} + \frac{1}{2} \right)
= \frac{1}{2}.
\]
In particular, $P(A) = \frac{1}{2} > 0$, so conditioning on $A$ is well-defined.

\medskip
\textbf{Step 2: Compute $P(X \le t,\, A)$ for all $t \in \mathbb{R}$.}
\begin{itemize}
    \item \emph{Case 1: $t \le 0$.}
    Since $X > 0$ with probability $1$, we have $0 \le P(X \le t,\, A) \le P(X \le t) = 0$, and hence
    \[
    P(X \le t,\, A) = 0.
    \]

    \item \emph{Case 2: $t > 0$.}
    Integrating the joint PDF over $\{(x,u) : 0 < x \le t,\ 0 < u < e^{-\lambda x}\}$, we get
    \begin{align*}
    P(X \le t,\, A)
    &= \int_{0}^{t} \int_{0}^{e^{-\lambda x}} \lambda e^{-\lambda x} \dd u \dd x \\
    &= \int_{0}^{t} \lambda e^{-2\lambda x} \dd x \\
    &= \left[ -\frac{1}{2} e^{-2\lambda x} \right]_{0}^{t} \\
    &= -\frac{1}{2} e^{-2\lambda t} - \left(-\frac{1}{2} e^{0}\right)
     = \frac{1}{2}\left(1 - e^{-2\lambda t}\right).
    \end{align*}
\end{itemize}

\medskip
\textbf{Step 3: Compute $P(X \le t \mid A)$ and identify the distribution.}
By the definition of conditional probability,
\[
P(X \le t \mid A) = \frac{P(X \le t,\, A)}{P(A)}.
\]
Substituting $P(A) = \frac{1}{2}$ and the two cases from Step 2:
\begin{itemize}
    \item For $t \le 0$:
    \[
    P(X \le t \mid A) = \frac{0}{1/2} = 0.
    \]
    \item For $t > 0$:
    \[
    P(X \le t \mid A)
    = \frac{\tfrac{1}{2}\left(1 - e^{-2\lambda t}\right)}{\tfrac{1}{2}}
    = 1 - e^{-2\lambda t}.
    \]
\end{itemize}
Combining both cases, for all real $t$:
\[
\boxed{
P(X \le t \mid A) =
\begin{cases}
0, & t \le 0, \\[4pt]
1 - e^{-2\lambda t}, & t > 0.
\end{cases}
}
\]
This is precisely the CDF of the Exponential distribution with rate parameter $2\lambda$. Therefore,
\[
\boxed{X \mid A \sim \mathrm{Expo}(2\lambda)}\,.
\]

\end{homeworkProblem}

\newpage

% =========================================================
% Problem 4
% =========================================================

\begin{homeworkProblem}{15}

Let $U_1,U_2,\ldots,U_n$ be i.i.d.
$\mathrm{Unif}(0,1)$ random variables.
Define

\[
Y_i=U_i(1-U_i),
\qquad
\overline Y_n=\frac1n\sum_{i=1}^{n}Y_i.
\]

\begin{enumerate}[(a)]

\item Find $\E[Y_i]$ and $\Var(Y_i)$.

\item Use the Central Limit Theorem to
find the smallest positive integer $n$
such that the approximate value of

\[
P\left(
\left|\overline Y_n-\frac16\right|
<\frac1{60}
\right)
\]

is at least $2\Phi(2)-1$, where $\Phi$
is the standard Normal CDF.

\end{enumerate}

\solution

\begin{enumerate}[(a)]

\item For any integer $k \ge 0$, the $k$-th moment of $U_i \sim \mathrm{Unif}(0,1)$ is
\[
\E[U_i^k] = \int_{0}^{1} u^k \dd u = \left[ \frac{u^{k+1}}{k+1} \right]_{0}^{1} = \frac{1}{k+1}.
\]
Using linearity of expectation (equivalently, LOTUS directly on $Y_i = U_i(1-U_i) = U_i - U_i^2$), we compute the mean of $Y_i$:
\begin{align*}
\E[Y_i]
&= \E[U_i - U_i^2]
 = \E[U_i] - \E[U_i^2] \\
&= \int_{0}^{1} u \dd u - \int_{0}^{1} u^2 \dd u \\
&= \frac{1}{2} - \frac{1}{3}
 = \frac{1}{6}.
\end{align*}
Next, expand $Y_i^2 = U_i^2(1-U_i)^2 = U_i^2(1 - 2U_i + U_i^2) = U_i^2 - 2U_i^3 + U_i^4$. Taking expectations gives the second moment:
\begin{align*}
\E[Y_i^2]
&= \E[U_i^2] - 2\E[U_i^3] + \E[U_i^4] \\
&= \frac{1}{3} - 2 \cdot \frac{1}{4} + \frac{1}{5} \\
&= \frac{1}{3} - \frac{1}{2} + \frac{1}{5}
 = \frac{10 - 15 + 6}{30}
 = \frac{1}{30}.
\end{align*}
Therefore, the variance of $Y_i$ is
\[
\Var(Y_i) = \E[Y_i^2] - \left(\E[Y_i]\right)^2
= \frac{1}{30} - \left(\frac{1}{6}\right)^2
= \frac{1}{30} - \frac{1}{36}
= \frac{6 - 5}{180}
= \frac{1}{180}.
\]
In summary,
\[
\boxed{\E[Y_i] = \frac{1}{6}, \qquad \Var(Y_i) = \frac{1}{180}}\,.
\]

\item Since $U_1, U_2, \ldots, U_n$ are i.i.d., the transformed random variables $Y_1, Y_2, \ldots, Y_n$ are also i.i.d.\ with common mean $\mu = \frac{1}{6}$ and common finite positive variance $\sigma^2 = \frac{1}{180}$. Consequently, the sample mean $\overline{Y}_n$ has mean and standard deviation
\[
\E[\overline{Y}_n] = \mu = \frac{1}{6},
\qquad
\mathrm{SD}(\overline{Y}_n) = \frac{\sigma}{\sqrt{n}} = \frac{1}{\sqrt{180n}}.
\]
Define the standardized sample mean
\[
Z_n = \frac{\overline{Y}_n - \mu}{\sigma / \sqrt{n}}
    = \sqrt{180n}\left(\overline{Y}_n - \frac{1}{6}\right).
\]
By the Central Limit Theorem, for large $n$ the distribution of $Z_n$ is approximated by a standard Normal random variable $Z \sim \mathcal{N}(0,1)$. Rewriting the target event in terms of $Z_n$, we obtain
\begin{align*}
P\left(\left|\overline{Y}_n - \frac{1}{6}\right| < \frac{1}{60}\right)
&= P\left(\left|\frac{\overline{Y}_n - \frac{1}{6}}{1/\sqrt{180n}}\right| < \frac{1/60}{1/\sqrt{180n}}\right) \\
&= P\left(|Z_n| < \frac{\sqrt{180n}}{60}\right).
\end{align*}
Simplifying the threshold on the right-hand side:
\[
\frac{\sqrt{180n}}{60}
= \frac{\sqrt{180n}}{\sqrt{3600}}
= \sqrt{\frac{180n}{3600}}
= \sqrt{\frac{n}{20}}.
\]
Applying the Central Limit Theorem approximation $Z_n \approx Z \sim \mathcal{N}(0,1)$ and using the symmetry property $\Phi(-z) = 1 - \Phi(z)$ of the standard Normal CDF, we get
\begin{align*}
P\left(\left|\overline{Y}_n - \frac{1}{6}\right| < \frac{1}{60}\right)
&\approx P\left(|Z| < \sqrt{\frac{n}{20}}\right) \\
&= \Phi\!\left(\sqrt{\frac{n}{20}}\right) - \Phi\!\left(-\sqrt{\frac{n}{20}}\right) \\
&= \Phi\!\left(\sqrt{\frac{n}{20}}\right) - \left(1 - \Phi\!\left(\sqrt{\frac{n}{20}}\right)\right) \\
&= 2\Phi\!\left(\sqrt{\frac{n}{20}}\right) - 1.
\end{align*}
We require this approximate probability to be at least $2\Phi(2) - 1$:
\[
2\Phi\!\left(\sqrt{\frac{n}{20}}\right) - 1 \ge 2\Phi(2) - 1
\iff
\Phi\!\left(\sqrt{\frac{n}{20}}\right) \ge \Phi(2).
\]
Because the standard Normal CDF $\Phi$ is strictly increasing on $\mathbb{R}$, this inequality holds if and only if
\[
\sqrt{\frac{n}{20}} \ge 2
\iff
\frac{n}{20} \ge 4
\iff
n \ge 80.
\]
Hence, the smallest positive integer $n$ satisfying the requirement is
\[
\boxed{n = 80}\,.
\]

\end{enumerate}

\end{homeworkProblem}

\newpage

% =========================================================
% Problem 5
% =========================================================

\begin{homeworkProblem}{20}

Let $B\sim\mathrm{Bin}(2,p)$, where
$0<p<1$.

Let $U_1,U_2,U_3$ be i.i.d.
$\mathrm{Unif}(0,1)$ random variables,
independent of $B$.

Define

\[
N=1+B,
\qquad
S=\sum_{i=1}^{N}U_i.
\]

\begin{enumerate}[(a)]

\item Find the moment generating function
$M_S(t)$ for all real $t$, including $t=0$.

\item Use the derivatives of $M_S(t)$
at $t=0$ to find $\E[S]$ and $\Var(S)$.

\end{enumerate}

\solution

\begin{enumerate}[(a)]

\item \textbf{Step 1: MGF of a single uniform random variable $U_i$.}
Let $M_U(t) = \E[e^{t U_1}]$ denote the common MGF of $U_1, U_2, U_3 \sim \mathrm{Unif}(0,1)$.
\begin{itemize}
    \item When $t = 0$:
    \[
    M_U(0) = \E[e^{0}] = \int_{0}^{1} 1\dd u = 1.
    \]
    \item When $t \neq 0$:
    \[
    M_U(t) = \int_{0}^{1} e^{tu}\dd u = \left[ \frac{e^{tu}}{t} \right]_{0}^{1} = \frac{e^t - 1}{t}.
    \]
\end{itemize}
By L'H\^opital's rule, $(e^t - 1)/t \to 1 = M_U(0)$ as $t \to 0$, so $M_U$ is continuous on all of $\mathbb{R}$.

\medskip
\textbf{Step 2: Conditional MGF of $S$ given $B$.}
Since $B \sim \mathrm{Bin}(2,p)$, its PMF is
\[
P(B = k) = \binom{2}{k} p^k (1-p)^{2-k}, \qquad k \in \{0, 1, 2\}.
\]
Conditioned on the event $\{B = k\}$, we have $N = 1 + k$, so $S = U_1 + \cdots + U_{k+1}$. Because $U_1, U_2, U_3$ are independent of $B$ and mutually i.i.d.,
\begin{align*}
\E\!\left[e^{tS} \mid B = k\right]
&= \E\!\left[\exp\!\left(t \sum_{i=1}^{k+1} U_i\right) \;\middle|\; B = k\right] \\
&= \E\!\left[\prod_{i=1}^{k+1} e^{t U_i}\right]
 = \prod_{i=1}^{k+1} \E\!\left[e^{t U_i}\right]
 = \left(M_U(t)\right)^{k+1}.
\end{align*}

\medskip
\textbf{Step 3: Apply the Law of Total Expectation.}
Summing over $k = 0, 1, 2$ and applying the Binomial Theorem,
\begin{align*}
M_S(t) = \E\!\left[e^{tS}\right]
&= \sum_{k=0}^{2} \E\!\left[e^{tS} \mid B = k\right] P(B = k) \\
&= \sum_{k=0}^{2} \left(M_U(t)\right)^{k+1} \binom{2}{k} p^k (1-p)^{2-k} \\
&= M_U(t) \sum_{k=0}^{2} \binom{2}{k} \left(p M_U(t)\right)^k (1-p)^{2-k} \\
&= M_U(t) \left(1 - p + p M_U(t)\right)^2.
\end{align*}
Substituting the explicit formula for $M_U(t)$ gives $M_S(t)$ for all $t \in \mathbb{R}$:
\[
\boxed{
M_S(t) =
\begin{cases}
\dfrac{e^t - 1}{t}\left(1 - p + p\,\dfrac{e^t - 1}{t}\right)^{\!2}, & t \neq 0, \\[12pt]
1, & t = 0.
\end{cases}
}
\]

\item \textbf{Step 1: Compute $M_U(0)$, $M_U'(0)$, and $M_U''(0)$.}
Using the Maclaurin series of $e^t$ (which converges for all $t \in \mathbb{R}$), we obtain for all $t \in \mathbb{R}$:
\[
M_U(t)
= \frac{e^t - 1}{t}
= \frac{1}{t}\sum_{k=1}^{\infty} \frac{t^k}{k!}
= \sum_{k=0}^{\infty} \frac{t^k}{(k+1)!}
= 1 + \frac{1}{2} t + \frac{1}{6} t^2 + \frac{1}{24} t^3 + \cdots,
\]
where the power series agrees with $M_U(0) = 1$ at $t = 0$. Differentiating term-by-term and evaluating at $t = 0$:
\begin{align*}
M_U'(t) &= \frac{1}{2} + \frac{1}{3} t + \frac{1}{8} t^2 + \cdots
\implies M_U'(0) = \frac{1}{2}, \\[4pt]
M_U''(t) &= \frac{1}{3} + \frac{1}{4} t + \cdots
\implies M_U''(0) = \frac{1}{3},
\end{align*}
along with $M_U(0) = 1$.

\medskip
\textbf{Step 2: Differentiate $M_S(t)$ via the chain and product rules.}
From part (a), $M_S(t) = g(M_U(t))$, where $g(m) = m(1 - p + pm)^2$. Differentiating $g(m)$ with respect to $m$:
\begin{align*}
g'(m)
&= (1 - p + pm)^2 + m \cdot 2(1 - p + pm) \cdot p
 = (1 - p + pm)(1 - p + 3pm), \\[4pt]
g''(m)
&= p(1 - p + 3pm) + 3p(1 - p + pm)
 = 4p(1 - p) + 6p^2 m.
\end{align*}
Evaluating $g'$ and $g''$ at $m = M_U(0) = 1$:
\[
g'(1) = (1 - p + p)(1 - p + 3p) = 1 + 2p,
\qquad
g''(1) = 4p(1 - p) + 6p^2 = 4p + 2p^2.
\]

\medskip
\textbf{Step 3: Evaluate $M_S'(0)$ and $M_S''(0)$ to find $\E[S]$ and $\Var(S)$.}
By the chain rule,
\[
M_S'(t) = g'(M_U(t))\, M_U'(t).
\]
Evaluating at $t = 0$ yields the expectation of $S$:
\[
\E[S] = M_S'(0) = g'(M_U(0))\, M_U'(0) = (1 + 2p) \cdot \frac{1}{2} = \frac{1}{2} + p.
\]
Differentiating $M_S'(t)$ once more using the product rule and the chain rule gives
\[
M_S''(t) = g''(M_U(t))\left(M_U'(t)\right)^2 + g'(M_U(t))\, M_U''(t).
\]
Evaluating at $t = 0$ yields the second moment of $S$:
\begin{align*}
\E[S^2] = M_S''(0)
&= g''(1)\left(M_U'(0)\right)^2 + g'(1)\, M_U''(0) \\*
&= (4p + 2p^2)\left(\frac{1}{2}\right)^2 + (1 + 2p) \cdot \frac{1}{3} \\
&= \left(p + \frac{1}{2} p^2\right) + \left(\frac{1}{3} + \frac{2}{3} p\right)
 = \frac{1}{3} + \frac{5}{3} p + \frac{1}{2} p^2.
\end{align*}
\newpage
Finally, we obtain the variance of $S$:
\begin{align*}
\Var(S)
&= \E[S^2] - \left(\E[S]\right)^2
 = M_S''(0) - \left(M_S'(0)\right)^2 \\
&= \left(\frac{1}{3} + \frac{5}{3} p + \frac{1}{2} p^2\right) - \left(\frac{1}{2} + p\right)^2 \\
&= \left(\frac{1}{3} + \frac{5}{3} p + \frac{1}{2} p^2\right) - \left(\frac{1}{4} + p + p^2\right) \\
&= \left(\frac{1}{3} - \frac{1}{4}\right) + \left(\frac{5}{3} - 1\right) p + \left(\frac{1}{2} - 1\right) p^2 \\
&= \frac{1}{12} + \frac{2}{3} p - \frac{1}{2} p^2
 = \frac{1 + 8p - 6p^2}{12}.
\end{align*}
Therefore,
\[
\boxed{
\E[S] = \frac{1}{2} + p,
\qquad
\Var(S) = \frac{1}{12} + \frac{2}{3} p - \frac{1}{2} p^2
        = \frac{1 + 8p - 6p^2}{12}
}\,.
\]

\end{enumerate}

\end{homeworkProblem}

\newpage

% =========================================================
% Problem 6
% =========================================================

\begin{homeworkProblem}{20}

Let $X$ be a positive continuous
random variable with PDF $f$, CDF $F$,
and failure (hazard) rate function $r$.

Suppose that

\[
r(t)=\frac{2t}{1+t^2},
\qquad t>0.
\]

Using the definitions of the failure
(hazard) rate function and the survival
function, prove that

\[
\boxed{\E[X]=\frac{\pi}{2}}.
\]

\solution

\begin{proof}
Let $S(t)$ denote the survival function of $X$, defined for $t \in \mathbb{R}$ by
\[
S(t) = P(X > t) = 1 - F(t).
\]
Since $X$ is a positive continuous random variable ($P(X > 0) = 1$), we have $F(0) = 0$ and thus
\[
S(0) = 1 - F(0) = 1.
\]
Furthermore, for $t > 0$, differentiating $S(t) = 1 - F(t)$ yields
\[
S'(t) = -F'(t) = -f(t).
\]

\medskip
\textbf{Step 1: Recover the survival function $S(x)$ from the hazard rate $r(t)$.}
By definition, the failure (hazard) rate function of $X$ is
\[
r(t) = \frac{f(t)}{1 - F(t)} = \frac{f(t)}{S(t)}, \qquad t > 0.
\]
Substituting $f(t) = -S'(t)$ gives the differential relation
\[
r(t) = -\frac{S'(t)}{S(t)} = -\frac{\diff}{\diff t}\ln S(t), \qquad t > 0.
\]
For any fixed $x > 0$, integrating both sides from $t = 0$ to $t = x$ and applying the Fundamental Theorem of Calculus yields
\[
\int_{0}^{x} r(t)\dd t
= -\int_{0}^{x} \frac{\diff}{\diff t}\ln S(t)\dd t
= -\Big[ \ln S(t) \Big]_{0}^{x}
= -\ln S(x) + \ln S(0).
\]
Since $S(0) = 1$, we have $\ln S(0) = 0$, and therefore
\[
-\ln S(x) = \int_{0}^{x} r(t)\dd t = \int_{0}^{x} \frac{2t}{1 + t^2}\dd t.
\]
To evaluate this integral, make the substitution $u = 1 + t^2$, so that $\diff u = 2t\dd t$. As $t$ increases from $0$ to $x$, $u$ increases from $1$ to $1 + x^2$. Hence,
\[
\int_{0}^{x} \frac{2t}{1 + t^2}\dd t
= \int_{1}^{1+x^2} \frac{1}{u}\dd u
= \Big[ \ln u \Big]_{1}^{1+x^2}
= \ln(1 + x^2) - \ln 1
= \ln(1 + x^2).
\]
It follows that
\[
-\ln S(x) = \ln(1 + x^2)
\implies
\ln S(x) = -\ln(1 + x^2) = \ln\!\left(\frac{1}{1 + x^2}\right).
\]
Exponentiating both sides, we obtain the survival function for all $x \ge 0$ (noting that $S(0) = 1$ also agrees with the formula):
\[
S(x) = \frac{1}{1 + x^2}, \qquad x \ge 0.
\]

\medskip
\textbf{Step 2: Express $\E[X]$ in terms of $S(t)$ and evaluate the integral.}
Since $X > 0$ with probability $1$ and $f(x) \ge 0$, writing $x$ as an integral over $(0,x)$ and interchanging the order of integration by Tonelli's theorem gives
\begin{align*}
\E[X]
&= \int_{0}^{\infty} x f(x)\dd x
 = \int_{0}^{\infty} \left( \int_{0}^{x} 1\dd t \right) f(x)\dd x \\
&= \int_{0}^{\infty} \left( \int_{t}^{\infty} f(x)\dd x \right) \dd t
 = \int_{0}^{\infty} P(X > t)\dd t
 = \int_{0}^{\infty} S(t)\dd t.
\end{align*}
Substituting $S(t) = \frac{1}{1 + t^2}$ into this identity and evaluating the improper integral, we obtain
\begin{align*}
\E[X]
&= \int_{0}^{\infty} \frac{1}{1 + t^2}\dd t
 = \lim_{b \to \infty} \int_{0}^{b} \frac{1}{1 + t^2}\dd t \\
&= \lim_{b \to \infty} \Big[ \arctan t \Big]_{0}^{b}
 = \lim_{b \to \infty} \arctan b - \arctan 0 \\
&= \frac{\pi}{2} - 0
 = \frac{\pi}{2}.
\end{align*}
Thus, $\E[X] = \frac{\pi}{2}$, completing the proof.
\end{proof}

\end{homeworkProblem}

\end{document}
